\begin{proof}
\textbf{Young’s Inequality.}
Let $a,b\ge0$ and let $p,q>1$ satisfy $\displaystyle\frac1p+\frac1q=1$. Then
\[
ab\le\frac{a^{p}}{p}+\frac{b^{q}}{q},
\]
with equality iff $a^{p}=b^{q}$ (including the trivial case $a=b=0$).

\medskip
\textbf{1. Degenerate cases.}
If $a=0$ or $b=0$, the left–hand side is $0$ while the right–hand side is a sum of non‑negative terms; thus the inequality holds, and equality occurs only when $a=b=0$. Hence we may assume $a>0$ and $b>0$.

\medskip
\textbf{2. Weighted AM–GM.}
Set
\[
\alpha=\frac1p,\qquad \beta=\frac1q,
\]
so that $\alpha,\beta>0$ and $\alpha+\beta=1$ because $\frac1p+\frac1q=1$.

The weighted arithmetic–geometric mean inequality states that for any $X,Y\ge0$,
\begin{equation*}
\alpha X+\beta Y\ge X^{\alpha}Y^{\beta},
\tag{2.1}
\end{equation*}
with equality iff $X=Y$.

Choose $X=a^{p}$ and $Y=b^{q}$. Then
\[
X^{\alpha}=(a^{p})^{1/p}=a,\qquad
Y^{\beta}=(b^{q})^{1/q}=b,
\]
and (2.1) becomes
\begin{equation*}
\frac{a^{p}}{p}+\frac{b^{q}}{q}\ge ab .
\tag{2.2}
\end{equation*}
Inequality (2.2) is precisely Young’s inequality.

\medskip
\textbf{3. Equality condition.}
Equality in (2.1) holds exactly when $X=Y$, i.e. when $a^{p}=b^{q}$. Under this relation the substitutions above give $\frac{a^{p}}{p}+\frac{b^{q}}{q}=ab$, so equality is attained in Young’s inequality. Conversely, if equality holds in Young’s inequality, then (2.2) must be an equality in the weighted AM–GM step, hence $a^{p}=b^{q}$. The trivial case $a=b=0$ also satisfies $a^{p}=b^{q}$.

\medskip
\textbf{4. Conclusion.}
For all non‑negative $a,b$ and conjugate exponents $p,q>1$,
\[
ab\le\frac{a^{p}}{p}+\frac{b^{q}}{q},
\]
with equality iff $a^{p}=b^{q}$ (including $a=b=0$). \qed
\end{proof}